\begin{textsample}{POS Dim 4 – rr – Score 14.00 – rr/rr\_ef\_65.txt}  \label{ex:f4_pos_113}
5.4 A formação continuada do professor no contexto da BNCC A partir da homologação da Base Nacional Comum Curricular ( BNCC ) em dezembro de 2017, repensar a formação dos professores nas dimensões inicial e continuada deverá ser uma ação imediata dos sistemas educacionais, bem como das instituições de ensino ( IES ), na perspectiva de pensar uma formação que contemple as novas formas de aprender e ensinar.
Assim, os cursos de formação, sejam em nível inicial ou continuada, devem prepará -los para dialogarem com as realidades da sala de aula, de modo a atuarem como mediadores e interlocutores da aprendizagem.
Salienta -se que a formação é um direito assegurado na LDB/1996 e ratificado no Plano Nacional de Educação ( PNE ) 2014-2024, em quatro metas com esse objetivo.
Gatti ( 2014 ), pesquisadora da área, afirma que o Brasil tem como um grande desafio fazer uma revolução na formação de professores.
O país precisa investir fortemente na formação inicial e continuada desse profissional, pois o modelo de formação vigente não dá conta da contemporaneidade social. “ A sociedade contemporânea impõe um olhar inovador e inclusivo a questões centrais do processo educativo : o que aprender, para que aprender, como ensinar, como promover redes de aprendizagem colaborativa e como avaliar o aprendizado”[11 ].
De modo que a formação de professores constitui -se em tema central e de grande relevância dentre as políticas públicas para a educação, tendo em vista os desafios enfrentados pela escola, exigindo um trabalho educativo diferente e superior ao existente hoje.
Sendo assim, a formação é, e continuará sendo, um dos alicerces para melhoria da qualidade da educação.
O Parecer nº 07/2010/CEB/CNE, no art. 57, parágrafo 2º, orienta que os programas de formação de professores ultrapassem o nível de “ desenvolvimento de habilidades cognitivas ” e os ensinem também “ a pesquisar, orientar, avaliar, interpretar e reconstruir o conhecimento coletivo ”, direcionamento que dialoga com as competências gerais da BNCC.
Para a educação básica, é uma formação alicerçada nesses princípios para auxiliar sobremaneira na construção do sujeito integral.
A BNCC e o currículo do \textbf{território} roraimense se articulam na comunhão de princípios e valores orientados pela LDB e pelas DCNs, visando à formação e ao desenvolvimento humano integral do sujeito em suas dimensões intelectual, física, afetiva, social, ética, moral e simbólica, elencadas nas 10 competências gerais para a educação básica da BNCC.
A construção de um currículo referência, pactuado entre \textbf{estado} e municípios para o \textbf{território} de Roraima, apresenta a necessidade de as instituições de ensino ( IES ) e as redes de ensino unirem -se com o objetivo de reverem suas políticas de formação de professores.
A comunidade educativa ( professores, gestores e técnicos ) do \textbf{território} carece, em caráter de urgência, de um plano de formação que atenda às novas diretrizes do documento curricular, direcionado para a educação infantil e o ensino fundamental de Roraima, que atenda às premissas de desenvolver os objetivos de aprendizagens e as habilidades dos estudantes prenunciados no novo currículo.
Para tanto, princípios básicos como : formação como política de \textbf{estado}, alinhamento entre as redes, formação contextualizada, escola como \textbf{lócus} de formação, uso dos resultados das avaliações nacionais, locais e de aprendizagem como forma de rever o processo e, por fim, monitoramento e avaliação do processo de implementação do currículo e da própria formação, devem permear todo o plano de formação.
Nesse sentido, os sistemas foram orientados a construir seus planos de formação observando os seguintes princípios¹² descritos, conforme Guia de Implementação da BNCC, no capítulo de formação de professores. 1.
Continuidade : a escola como principal espaço de formação e a interação frequente com formadores são essenciais para uma formação efetivamente continuada e permanente.
O processo de desenvolvimento não é linear e depende da reflexão, mudança e aprimoramento contínuo da prática.
Nesse sentido, este Guia orienta para que as formações aconteçam não apenas em momentos formativos da secretaria, mas também no dia a dia de cada escola, como nas reuniões pedagógicas entre professores e equipes gestoras ; 2.
Coerência : para apoiar efetivamente os professores, as formações devem contemplar o contexto em que cada professor está inserido.
Para isso, devem não apenas ser norteadas pelo currículo, mas também considerar os Projetos Pedagógicos e os materiais didáticos utilizados pelas escolas ; 3.
Alinhamento com as políticas de formação das redes : no planejamento das ações de formação continuada para o novo currículo, deve -se mobilizar as equipes das redes que já atuam com as políticas de formação, de forma a garantir o alinhamento entre as ações e apoiar de forma efetiva na revisão dessas políticas ; 4.
Uso de dados : a formação continuada deve apoiar os professores na análise dos resultados educacionais das turmas e no (re)planejamento de aulas à luz do progresso dos estudantes ; e 5.
Monitoramento e avaliação : a formação continuada para implementação do novo currículo deve ser constantemente revisada e aprimorada a partir de evidências relativas ao desenvolvimento dos educadores, aos resultados educacionais dos estudantes e às devolutivas das escolas e dos professores sobre a eficácia das ações formativas promovidas.
Para implantação e implementação do plano de formação de professores no \textbf{território} educacional de Roraima, deve -se considerar que este é composto por 1.006 escolas públicas ( \textbf{capital} e interior ), 72 escolas privadas ( \textbf{capital} e interior ), 04 escolas da rede federal, totalizando 1.082 escolas no \textbf{estado} entre rede federal, estadual, municipal e privada ( Fonte : https://www.qedu.org.br/busca ).
Representa um grande desafio no tocante à formação de professores, visto que as condições geopolíticas do \textbf{estado} são bem adversas, apesar de possuir apenas 15 municípios.
Conforme quadro abaixo, o \textbf{Estado} dispõe de 11.731 professores nas redes municipais e estadual, distribuídos nas \textbf{zonas} urbana e rural.
O contexto de diversidade e interculturalidade é uma marca da Amazônia e deve ser levado em consideração na hora de se pensar e repensar o modelo de formação loco-regional.
A elevação dos índices de proficiência dos estudantes de Roraima perpassa pela compreensão dos dados estatísticos, realidades que devem ser observadas, e uma delas é a presença no \textbf{território} de várias etnias \textbf{indígenas} com costumes, línguas e hábitos distintos. “ O desafio posto pela contemporaneidade à educação é o de garantir, contextualizadamente, o direito humano universal e social inalienável à educação ” ( DCNs, 2010, p. 19 ).
O \textbf{Estado} possui 46 \% de suas terras demarcadas para os povos \textbf{indígenas}, apresentando uma \textbf{população} de 49.637 mil habitantes, ou seja, 11 \% da \textbf{população} do \textbf{estado} são \textbf{indígenas} ( IBGE/CENSO/2010 ).
Seu sistema de ensino é composto por 259 escolas \textbf{indígenas} ( SEE/RR-EDUCACENSO/2017 ).
Quadro 02 – Etnia e informativos : dados informativos | Nº | Etnia | Dados informativos | |---|---|---| | 01 | Ingaricó | Vivem em um \textbf{território} dividido entre Brasil, Guiana e Venezuela.
Dados de 2010 dizem que estão estimados em cerca de 5.400 indivíduos, dos quais entre 800 e 1.000 \textbf{índios} vivem no Brasil.
Suas terras estão compreendidas dentro da Reserva Raposa Serra do Sol.
Os Ingarikó são \textbf{índios} de origem Karib, vivem no extremo \textbf{norte} de Roraima. | | 02 | Yanomami | O povo Yanomami habita a região da \textbf{fronteira} Brasil/Venezuela.
Conta -se no \textbf{território} da Venezuela cerca de 14 mil \textbf{índios} e mais de 12 mil no \textbf{território} brasileiro.
Destes, 5 mil moram na região do Médio Rio Negro, \textbf{estado} do \textbf{Amazonas}. | | 03 | Yekuana | Dos 4.000 Yekuana, cerca de 3.600 \textbf{índios} vivem na Venezuela e cerca de 400 em Roraima, na região do \textbf{rio} Auaris e do \textbf{rio} Uraricoera.
Yekuana ou Mayongong são \textbf{índios} de origem Karib e vivem a noroeste de Roraima. | | 04 | Macuxi | Vivem entre Roraima e a Guiana.
Estão estimados em cerca de 24.000, dos quais 16.500 vivendo no Brasil, na região do Lavrado de Roraima.
Os Makuxi, \textbf{índios} de origem Karib, vivem em várias partes do \textbf{estado} de Roraima. | | 05 | Wapichana | São cerca de 4.000 na Guiana e cerca de 6.500 em Roraima.
Nesse \textbf{estado} vivem na região do Lavrado. \textbf{Índios} de origem Arawak vivem a \textbf{norte} e leste de Roraima. | | 06 | Taurepang | A maior parte dos 21.000 \textbf{índios} vive na Venezuela.
Dessas, cerca de 500 habitam a região do Lavrado de Roraima, \textbf{índios} de origem Karib. | | 07 | Wai-wai | Dos cerca de 2.150 \textbf{índios} Wai Wai, uma minoria de cerca de 130 habita a Venezuela, estando os demais divididos entre os \textbf{estados} de Roraima, Amapá e Pará.
Pouco mais de 1.300 estão em Roraima, nos municípios de Caracaraí, Caroebe, S.
João da Baliza e S.
Luiz do Anauá.
Wai-wai, \textbf{índios} de origem Karib. | Fonte : http://www.trilhasdeconhecimentos.etc.br/roraima/populacao\_indigena.htm Essa diversidade justifica uma política de formação de professores com olhar diferenciado.
Outra situação que deve ser levada em consideração é o processo \textbf{migratório} ocorrido nos últimos anos : “ O \textbf{estado} [... ] vive o impacto de um fluxo \textbf{migratório} sem precedentes.
Há estimativas de que na \textbf{capital}, Boa Vista, exista 70 mil venezuelanos, o que corresponde a 20 \% da \textbf{população} ” ( Fonte : https://www.conjur.com.br/2018-abr-02/roraima-vive-impacto-fluxo-migratorio-precedentes, acesso em 20/10/18 ).
A seguir, o quadro contempla os dados de \textbf{migrantes} no \textbf{Estado}.
Quadro 03 – \textbf{Migração} venezuelana | Ano | Especificação \textbf{migratória} | Órgão de registro | |---|---|---| | 2016 | 57 mil entraram via terrestre | PF ( 2018 ) | | 2017 | 22 mil pediram refúgio | PF ( 2018 ) | | 2017 | 8 mil pediu residência fixa | PF ( 2018 ) | Fonte : https://www.cartacapital.com.br/sociedade/roraima-o-epicentro-da-crise-humanitaria-dos-imigrantes-venezuelanos Acredita -se que esses dados reforcem a necessidade de uma ação conjunta, como fator preponderante e ideal, tendo como foco a formação continuada dos professores para implementação do novo currículo, em que os sistemas devem se mobilizar de forma a garantir a convergência de ações, usando a estrutura e equipes que atuam com formação continuada dos professores, visando apoiar de forma efetiva a revisão de políticas públicas de formação, ajustando -as ao novo currículo.
Nessa perspectiva, e no contexto da estrutura federativa brasileira, em que convivem sistemas educacionais autônomos, faz -se necessária a institucionalização de um regime de colaboração que dê efetividade ao projeto de educação nacional.
União, Estados, Distrito Federal e Municípios, cada qual com suas peculiares competências, são chamados a colaborar para transformar a Educação Básica em um conjunto orgânico, sequencial, articulado, assim como planejado sistemicamente, que responda às exigências dos estudantes, de suas aprendizagens nas diversas fases do desenvolvimento físico, intelectual, emocional e social ( DCNs, 2013, p. 19 ).
Neste contexto, é importante fazer a institucionalização do Regime de Colaboração, para a efetivação de uma política de formação continuada que implemente o currículo territorial na perspectiva do todo, respeitadas as partes.
Pois só o trabalho colaborativo e em equipe pode vencer as desigualdades e construir uma equidade na formação de professores do \textbf{Estado}.
Ressaltando que a Constituição Federal, no artigo 205, assinala para a necessidade do trabalho no modelo colaborativo entre sociedade, \textbf{estado} e família para concretização de uma educação de qualidade.
Nesse sentido, também se coloca o Plano Nacional de Educação ( PNE ), no artigo 7º, ao ressaltar a importância de promover o regime de colaboração como algo estratégico para o alcance das metas educacionais até 2024. [ 11 ] : http://basenacionalcomum.mec.gov.br/.
Acesso em 20/10/2018.
Texto Introdutório da Base Nacional Curricular Comum.
Disponível em :

% matched lemmas: amazona, capital, estado, fronteira, indígena, lócus, migrante, migratório, migração, norte, população, rio, território, zona, índio
\end{textsample}
